\documentclass{article}
\usepackage{graphicx}
\usepackage{amsmath}
\usepackage{amsfonts}
\usepackage{bookmark}
\usepackage{booktabs}
\usepackage{exercises}
\graphicspath{ {} }
\usepackage{indentfirst}
\usepackage{amsthm}
\usepackage{yhmath}
\usepackage{geometry}
\usepackage{placeins}
\usepackage{csquotes}
\usepackage{pgfplots}
\newcommand*{\oldneg}{\mathord{\sim}}
\newcommand\showdiv[1]{\overline{\smash{)}#1}}
\newcommand\Mydiv[1]{%
\smash{\raise.3ex\hbox{$\big)$}}\mkern-8mu
        \overline{\enspace\strut#1}}
\setlength{\tabcolsep}{5pt}
\theoremstyle{definition}
\newtheorem{definition}{Definition}

\setlength{\parindent}{.5in}
\setlength{\parskip}{12pt}
\renewcommand{\arraystretch}{2}
\geometry{margin=1in}
\newcommand{\?}{\stackrel{?}{=}}
\pgfplotsset{width=7cm,compat=1.18}\usepgfplotslibrary{statistics}

\title{Chapter 10 Test (pg. 673)}
\begin{document}
\maketitle

\section*{1}
We are given the function $f(x)=4\cos 2x$. This has an amplitude of 4 and a period of $\frac{2\pi}{2}=\pi$. Now we identify the five key points: It has 
$x$-intercepts at $\frac{\pi}{4}$ and $\frac{3\pi}{4}$. It has maxima at $(0,4)$ and $(\pi,4)$ and minima at $(\frac{\pi}{2},-4)$:\newline

\section*{2}
We are given the function $y=\frac{3}{2}\sin\pi x$. This has an amplitue of $\frac{3}{2}$ and a period of $\frac{2\pi}{\pi}=2$. Now we identify the five key points: It has 
$x$-intercepts at $0$, $2$ and $\frac{2}{2}=1$. It has a maximum at $(\frac{2}{4},\frac{3}{2})=(\frac{1}{2},\frac{3}{2})$ and a minimum at $(\frac{3}{2},\frac{3}{2})$:\newline

\section*{3}
We are given the function $f(x)=-4\tan\frac{\pi}{2} x$. Ths has a period of $\frac{\pi}{\frac{\pi}{2}}=2$ and asymptotes at odd multiples of $\dfrac{\pi}{2\frac{\pi}{2}}=1$. Now we find the key points of the 
graph of $y=4\tan\frac{\pi}{2}x$. It has halfway points of $(-\frac{\pi}{4},-4)$ and $(\frac{\pi}{4},4)$. Thus the graph of the function 
$f(x)=-4\tan\frac{\pi}{2}x$ (which is a reflection of the other graph) has halfway points at $(-\frac{\pi}{4},4)$ and $(\frac{\pi}{4},-4)$. In addition, the graph has an $x$-intercept at 0:\newline

\section*{4}
We are given the function $y=\sin(x-\pi)-2$. This has a period of $2\pi$ and an amplitude of $1$. Also, it has a horizontal shift of $\pi$ and a vertical shift of $-2$, meaning the midline is $y=-2$. Now we find the five 
key points: The points on the midline are $0+\pi=\pi$, $\frac{1}{2}+\pi$, and $1+\pi$. The maximum is $(\frac{1}{4}\times 1+\pi,1-2)=(\frac{1}{4}+\pi,-2)$ and the minimum is
$(\frac{3}{4}\times 1+\pi,-1-2)=(\frac{3}{4}+\pi,-3)$:\newline

\section*{5}
We are given the function $f(x)=3\tan\left(x-\frac{\pi}{2}\right)$. It has a period of $\pi$ and horizontal asymptotes at odd multiples of 
$\frac{\pi}{2}+\frac{\pi}{2}=\pi$. Now we identify the key points of the graph: It has an $x$-intercept at $0+\frac{\pi}{2}=\frac{\pi}{2}$, and halfway points at
$(-\frac{\pi}{4}+\frac{\pi}{2},-3)=(\frac{\pi}{4},-3)$ and $(\frac{\pi}{4}+\frac{\pi}{2},3)=(\frac{3\pi}{4},3)$:\newline

\section*{6}
We are given the function $y=-2\cos\frac{1}{3}x+3$. It has an amplitude of $2$ and a period of $\frac{2\pi}{\frac{1}{3}}=6\pi$. Also, it has a vertical shift of $3$, meaning the midlein is $y=3$. Now we find the five 
key points of the graph of $y=2\cos\frac{1}{3}x+3$: The points on the midline are $\frac{1}{4}\times 6\pi=\frac{3\pi}{2}$ and $\frac{3}{4}\times 6\pi=\frac{9\pi}{2}$. The maxima are $(0,2+3)=(0,5)$ and 
$(6\pi,2+3)=(6\pi,5)$, while the minimum is $(\frac{6\pi}{2},-2+3)=(3\pi,1)$. Thus the graph of the function $y=-2\cos\frac{1}{3}x+3$ (which is a reflection of the other graph) has the same points on the midline, 
minima at $(0,1)$ and $(6\pi,1)$, and a maximum at $(3\pi,5)$:\newline

\section*{7}
\[ \frac{\sin(-\theta)}{\tan(-\theta)} \]
\[ =\frac{-\sin\theta}{-\tan\theta} \]
\[ =\frac{\sin\theta}{\tan\theta} \]
\[ =\cos\theta\textrm{ (Tangent Identity)} \]

\section*{8}
\[ \cos^2 x+\sin^2 x-\csc^2 x \]
\[ =1-\csc^2 x\textrm{ (Pythagorean Identity)} \]
\[ =-\cot^2 x\textrm{ (Pythagorean Identity)} \]

\section*{9}
\[ \dfrac{\sin\left(\frac{\pi}{2}-x\right)}{\sec x} \]
\[ =\frac{\cos x}{\sec x} \]
\[ =\dfrac{\cos x}{\frac{1}{\cos x}} \]
\[ =\cos^2 x \]

\section*{10}
We are given a maximum at $(2,5)$ and a minimum at $(1,-1)$. Thus the amplitude is $\frac{5+1}{2}=3$ and the vertical shift is $\frac{5-1}{2}=2$. Also, the period is $2(2-1)=2$ meaning
\[ 2=\frac{2\pi}{b} \] \[ b=\pi \] Because the $x$-coordinate of the 
minimum is half of the maximum, we can model this as a cosine function:
\[ y=3\cos\pi x+2 \]

\section*{11}
We are given a maximum at $(\frac{9\pi}{8},1)$ and a minimum at $(\frac{3\pi}{8},-5)$. Thus the amplitude is $\frac{1+5}{2}=3$ and the vertical shift is $\frac{1-5}{2}=-2$. Also, the period is 
$2(\frac{9\pi}{8}-\frac{3\pi}{8})=\frac{3\pi}{2}$ meaing \[ \frac{3\pi}{2}=\frac{2\pi}{b} \] \[ b=\frac{4}{3} \] Because the $x$-coordinate of the minimum is one third of the maximum we can model this as a reflection of 
a sine function with no horizontal shift:
\[ y=-3\sin\frac{4}{3}x-2 \]

\section*{12}
\[ \cos 3x=\cos(x+2x) \]
\[ =\cos x\cos 2x-\sin x\sin 2x \]
\[ =\cos x(2\cos^2 x-1)-2\sin^2 x\cos x \]
\[ =2\cos^3 x-\cos x-2\sin^2 x\cos x \]
\[ =2\cos^3 x-\cos x(1+2\sin^2 x) \]
\[ =2\cos^3 x-\cos x(1+2(1-\cos^2 x)) \]
\[ =2\cos^3 x-\cos x(3-2\cos^2 x) \]
\[ =2\cos^3 x-3\cos x+2\cos^3 x \]
\[ =4\cos^3 x-3\cos x \]

\section*{13}
\[ 9\sin^2 x\tan x=16\tan x \]
\[ 9\sin^2 x\tan x-16\tan x=0 \]
\[ \tan x(9\sin^2 x-16)=0 \]
\[ \tan x=0 \vee \sin x=\pm\frac{4}{3} \]
The second option is ruled out, as the domain of $\sin x$ is only from $-1$ to $1$. Thus $\tan x=0$, which, in the interval $0\le x\le 2\pi$, holds at $0$ and $\pi$. We check these solutions:
\[ 9\sin^2 0\tan 0\? 16\tan 0 \]
\[ 0=0 \]
and
\[ 9\sin^2 \pi\tan\pi\? 16\tan\pi \]
\[ 0=0 \] meaning neither of the solutions are extraneous.

\section*{14}
\[ (1-\tan^2 x)\tan 2x=2\sqrt{3} \]
\[ (1-\tan^2 x)\frac{2\tan x}{1-\tan^2 x}=2\sqrt{3} \]
\[ 2\tan x=2\sqrt{3} \]
\[ \tan x=\sqrt{3} \]
In the interval $0\le x\le 2\pi$, the solutions are $\frac{\pi}{3}$ and $\pi+\frac{\pi}{3}=\frac{4\pi}{3}$. We check these solutions:
\[ (1-\tan^2\frac{\pi}{3})\tan 2\frac{\pi}{3}\? 2\sqrt{3} \]
\[ (1-3)\tan \frac{2\pi}{3}\? 2\sqrt{3} \]
\[ -2(-\sqrt{3})\? 2\sqrt{3} \]
\[ 2\sqrt{3}=2\sqrt{3} \]
and
\[ (1-\tan^2\frac{4\pi}{3})\tan 2\frac{4\pi}{3}\? 2\sqrt{3} \]
\[ -2(-\sqrt{3})=2\sqrt{3} \]
meaning neither of teh solutions are extraneous.

\section*{15}
\[ \sin\frac{x}{2}=\frac{\sqrt{2}}{2} \]
\[ \pm\sqrt{\frac{1-\cos x}{2}}=\frac{\sqrt{2}}{2} \]
\[ \frac{1-\cos x}{2}=\frac{1}{2} \]
\[ 1-\cos x=1 \]
\[ \cos x=0 \]
In the interval $0\le x\le 2\pi$, the solutions are $\frac{\pi}{2}$ and $\frac{3\pi}{2}$. We check these solutions:
\[ \sin\dfrac{\frac{\pi}{2}}{2}\?\frac{\sqrt{2}}{2} \]
\[ \frac{1}{\sqrt{2}}\?\frac{\sqrt{2}}{2} \]
\[ \frac{\sqrt{2}}{2}=\frac{\sqrt{2}}{2} \]
and
\[ \sin\dfrac{\frac{3\pi}{2}}{2}\?\frac{\sqrt{2}}{2} \]
\[ \frac{1}{\sqrt{2}}\?\frac{\sqrt{2}}{2} \]
\[ \frac{\sqrt{2}}{2}=\frac{\sqrt{2}}{2} \]
meaning neither of the solutions are extraneous.

\section*{16}
\[ 6\tan^2x-2=0 \]
\[ \tan^2x=\frac{1}{3} \]
\[ \tan x=\pm\frac{1}{\sqrt{3}} \]
In the interval $0\le x\le\pi$, the solutions are $\frac{\pi}{6}$ and $\pi-\frac{\pi}{6}=\frac{5\pi}{6}$. We check these solutions:
\[ 6\tan^2\frac{\pi}{6}-2\? 0 \]
\[ 2-2=0 \]
and
\[ 6\tan^2\frac{5\pi}{6}-2\? 0 \]
\[ 2-2=0 \]
meaning neither of the solutions are extraneous. Hence the general solution is
\[ x=\frac{\pi}{6}+\pi n\textrm{ or }x=\frac{5\pi}{6}+\pi n \] where $n$ is an integer.

\section*{17}
\[ \cos x=\sin 2x\sin x \]
\[ \cos x=2\sin^2 x\cos x \]
\[ 2\sin^2 x\cos x-\cos x=0 \]
\[ \cos x(2\sin^2 x-1)=0 \]
\[ \cos x=0\vee \sin x=\pm\frac{1}{\sqrt{2}} \]
In the interval $0\le x\le 2\pi$, the solutions to the first option are $\frac{\pi}{2}$ and $\frac{3\pi}{2}$, while the solutions of the second are
\[ \frac{\pi}{4},\pi+\frac{\pi}{4}=\frac{5\pi}{4},\pi-\frac{\pi}{4}=\frac{3\pi}{4},2\pi-\frac{\pi}{4}=\frac{7\pi}{4} \]
We check all of these solutions:
\[ \cos\frac{\pi}{2}\?\sin 2\frac{\pi}{2}\sin\frac{\pi}{2} \]
\[ 0=0 \]
and
\[ \cos\frac{3\pi}{2}\?\sin 2\frac{3\pi}{2}\sin\frac{3\pi}{2} \]
\[ 0=0 \]
and
\[ \cos\frac{\pi}{4}\? \sin 2\frac{\pi}{4}\sin\frac{\pi}{4} \]
\[ \frac{1}{\sqrt{2}}=\frac{1}{\sqrt{2}} \]
and
\[ \cos\frac{3\pi}{4}\?\sin 2\frac{3\pi}{4}\sin\frac{3\pi}{4} \]
\[ -\frac{1}{\sqrt{2}}=-\frac{1}{\sqrt{2}} \]
and
\[ \cos\frac{5\pi}{4}\?\sin 2\frac{5\pi}{4}\sin\frac{5\pi}{4} \]
\[ -\frac{1}{\sqrt{2}}=-\frac{1}{\sqrt{2}} \]
and
\[ \cos\frac{7\pi}{4}\?\sin 2\frac{7\pi}{4}\sin\frac{7\pi}{4} \]
\[ \frac{1}{\sqrt{2}}=\frac{1}{\sqrt{2}} \]
meaning none of the solutions are extraneous. Hence the general solution is
\[ x=\frac{\pi}{2}+\pi n\textrm{ or }x=\frac{\pi}{4}+\pi n\textrm{ or }x=\frac{3\pi}{4}+\pi n \] where $n$ is an integer.

\section*{18}
\[ \sin\frac{x}{2}=1-\cos x \]
\[ \pm\sqrt{\frac{1-\cos x}{2}}=1-\cos x \]
\[ \frac{1-\cos x}{2}=(1-\cos x)^2 \]
\[ \frac{1-\cos x}{2}=1+\cos^2 x-2\cos x \]
\[ 1-\cos x=2+2\cos^2 x-4\cos x \]
\[ 2\cos^2 x-3\cos x+1=0 \]
\[ (2\cos x-1)(\cos x-1)=0 \]
\[ \cos x=\frac{1}{2}\textrm{ or }\cos x=1 \]
In the interval $0\le x\le 2\pi$, the solutions to the first option are $\frac{\pi}{3}$ and $2\pi-\frac{\pi}{3}=\frac{5\pi}{3}$ while the solution to the second is
$0$. We check these solutions:
\[ \sin\dfrac{\frac{\pi}{3}}{2}\? 1-\cos\frac{\pi}{3} \]
\[ \frac{1}{2}\? 1-\frac{1}{2} \]
\[ \frac{1}{2}=\frac{1}{2} \]
and
\[ \sin\dfrac{\frac{5\pi}{3}}{2}\? 1-\cos\frac{5\pi}{3} \]
\[ \frac{1}{2}=\frac{1}{2} \]
and
\[ \sin\dfrac{0}{2}\? 1-\cos 0 \]
\[ 0=0 \]
meaning none of the solutions are extraneous. Thus the general solution of the equation is
\[ x=2\pi n\textrm{ or }x=\frac{\pi}{3}+\pi n \] where $n$ is an integer.


\section*{19}
\[ \sin 255^\circ=\sin (180^\circ+(30^\circ+45^\circ)) \]
\[ =\sin 180^\circ\cos(30^\circ+45^\circ)+\sin (30^\circ+45^\circ)\cos 180^\circ \]
\[ =-(\sin 30^\circ\cos 45^\circ+\sin 45^\circ\cos 30^\circ) \]
\[ =-(\frac{1}{2}\frac{1}{\sqrt{2}}+\frac{1}{\sqrt{2}}\frac{\sqrt{3}}{2}) \]
\[ =-\frac{1+\sqrt{3}}{2\sqrt{2}} \]
\[ =-\frac{\sqrt{2}+\sqrt{6}}{4} \]

\section*{20}
Note that $-\frac{\pi}{8}$ is coterminal with $-\frac{\pi}{8}+2\pi=\frac{15\pi}{8}$ and lies in Quadrant IV.
\[ \cos(-\frac{\pi}{8})=\cos(-\dfrac{\frac{\pi}{4}}{2}) \]
\[ =\pm\sqrt{\frac{1+\cos(-\frac{\pi}{4})}{2}} \]
\[ =\sqrt{\frac{1+\frac{1}{\sqrt{2}}}{2}}\textrm{ (Because $-\frac{\pi}{8}$ lies in Quadrant IV)} \]
\[ =\sqrt{\dfrac{\frac{\sqrt{2}+1}{\sqrt{2}}}{2}} \]
\[ =\sqrt{\frac{\sqrt{2}+1}{2\sqrt{2}}} \]

\section*{21}
\[ \tan\frac{5\pi}{12}=\tan\dfrac{\frac{5\pi}{6}}{2} \]
\[ =\frac{\sin\frac{5\pi}{6}}{1+\cos\frac{5\pi}{6}} \]
\[ =\frac{\frac{1}{2}}{1-\frac{\sqrt{3}}{2}} \]
\[ =\dfrac{\frac{1}{2}}{\frac{2-\sqrt{3}}{2}} \]
\[ =\frac{1}{2-\sqrt{3}} \]
\[ =\frac{2+\sqrt{3}}{(2-\sqrt{3})(2+\sqrt{3})} \]
\[ =\frac{2+\sqrt{3}}{4-3} \]
\[ =2+\sqrt{3} \]

\section*{22}
\[ \sin\frac{10\pi}{3}=\sin(2\pi+\frac{4\pi}{3}) \]
\[ =\sin 2\pi\cos\frac{4\pi}{3}+\sin\frac{4\pi}{3}\cos 2\pi \]
\[ =\sin\frac{4\pi}{3} \]
\[ =-\frac{\sqrt{3}}{2} \]

\section*{23}
We are give a diameter of 11 feet, that there are 15 revolutions per minute, and a maximum of 2 feet below the water. Thus the minimum is $17$ feet below the wheel. We are writing a model that starts on a point 
at the very top of the wheel; hence the maximum is at the $y$-axis, meaning we can model this as a cosine function. Also, we know that the period is $\frac{1}{15}$, meaning
\[ \frac{1}{15}=\frac{2\pi}{b} \] \[ b=30\pi \] Also, the amplitude is $\frac{-2+17}{2}=7.5$ and the midline is $\frac{-2-17}{2}=-9.5$.
Thus the model is
\[ y=7.5\cos 30\pi x-9.5 \]

\section*{24}
I was unsure of how to preform this calculation on my calculator.

\end{document}
